\documentclass[11pt,a4paper]{article}

\usepackage{ctex}
\usepackage[T1]{fontenc}
\usepackage[top=2.5cm, bottom=2.5cm, left=2.5cm, right=2.5cm]{geometry}
\usepackage{booktabs}
\usepackage{array}
\usepackage{enumitem}
\usepackage{tabularx}
\usepackage{amssymb}
\usepackage{xcolor}
\definecolor{primary}{HTML}{1a1a2e}
\definecolor{accent}{HTML}{3b82f6}
\definecolor{success}{HTML}{22c55e}
\definecolor{danger}{HTML}{ef4444}
\definecolor{warning}{HTML}{f59e0b}
\definecolor{graytext}{HTML}{6b7280}
\usepackage[hidelinks]{hyperref}
\hypersetup{colorlinks=true, linkcolor=primary, urlcolor=accent}
\usepackage{titlesec}
\titleformat{\section}{\Large\bfseries\color{primary}}{}{0em}{}[\titlerule]
\titleformat{\subsection}{\large\bfseries\color{primary}}{}{0em}{}
\usepackage{fancyhdr}
\pagestyle{fancy}
\fancyhf{}
\fancyhead[L]{\small\color{graytext} 企业级 AI Agent 应用商店}
\fancyhead[R]{\small\color{graytext} MRD}
\fancyfoot[C]{\small\color{graytext} \thepage}
\renewcommand{\headrulewidth}{0.4pt}

\title{\textbf{MRD：市场需求文档}\\[0.5em]\large Market Requirements Document}
\author{万能产品小助手（PM AI）\\魏子政 审阅}
\date{2026-05-06 / 版本 v1.0}

\begin{document}

\maketitle

\begin{center}
\begin{tabular}{>{\bfseries}p{3cm} p{10cm}}
\toprule
文档类型 & MRD（Market Requirements Document） \\
定位 & 市场需求文档 -- 这东西有没有人要？市场有多大？竞品是什么？ \\
核心问题 & 先想清楚再动手 \\
\bottomrule
\end{tabular}
\end{center}

\vspace{1em}

\section{产品定位}

企业级 AI Agent 应用商店：面向企业内部的 AI Agent / Skill 统一管理平台。开发者提交 Skill 包，管理员审批上架，员工浏览、搜索、一键安装。平台提供三层安全扫描（Semgrep + ClamAV + LLM 语义分析）、邮箱认证注册、完整的应用生命周期管理（上架/下架/重新提交）。

\section{目标用户}

\begin{table}[h]
\centering
\begin{tabularx}{\linewidth}{lXX}
\toprule
\textbf{角色} & \textbf{是谁} & \textbf{核心诉求} \\
\midrule
管理员 & IT/运维部门 & 统一管控所有 Skill，审批上架/下架，安全合规 \\
开发者 & 内部各业务团队 & 提交 Skill 包，版本管理，查看审批状态 \\
普通用户 & 全体员工 & 发现好用的 Skill，搜索筛选，一键安装 \\
\bottomrule
\end{tabularx}
\end{table}

\section{痛点与市场机会}

\begin{enumerate}[nosep]
  \item 企业内部 AI Agent 越来越多，分散在各处，没有统一入口
  \item 缺乏安全审核机制，Skill 权限和行为不可控
  \item 员工不知道有哪些 Skill 可用，重复开发严重
  \item 没有 Skill 的生命周期管理（上架、下架、更新、卸载）
\end{enumerate}

\section{竞品分析}

\begin{table}[h]
\centering
\begin{tabularx}{\linewidth}{lXXX}
\toprule
\textbf{方案} & \textbf{类型} & \textbf{优势} & \textbf{不足} \\
\midrule
OpenClaw skills.sh & 开源 Skill 生态 & 56+ Agent，SKILL.md 标准，CLI 安装 & 无审批流，无安全扫描，非企业级 \\
Flow Nexus & 开源 Skill 商店 & MCP 架构，模板部署 & 不完整，无企业级审批/审计 \\
Membrane & 开源中间件 & 免密钥代理，40+ 模型 & 偏基础设施，非应用商店 \\
ChatGPT GPT Store & SaaS & 用户量大，生态成熟 & 不适合企业私有部署 \\
自建 & 定制开发 & 完全可控 & 需要投入开发资源 \\
\bottomrule
\end{tabularx}
\end{table}

\textbf{结论}：现有方案均不满足企业级需求（审批工作流 + 三层安全扫描 + 生命周期管理）。本项目需自研，借鉴 OpenClaw skills.sh 的 SKILL.md 标准和分类体系。

\section{核心价值}

\begin{itemize}[nosep]
  \item \textbf{安全合规}：三层安全扫描 + 管理员审批，降低 Skill 安全风险
  \item \textbf{统一入口}：一个平台管理所有 Skill，搜索 + 分类 + 筛选
  \item \textbf{生命周期}：完整的上架/下架/重新提交流转，Skill 状态全程可追踪
  \item \textbf{一键安装}：支持 zip 下载 + OpenClaw 本地安装，用户体验简单直接
  \item \textbf{邮箱认证}：注册需邮箱验证，提升账号安全等级
\end{itemize}

\section{协作模式}

\begin{itemize}[nosep]
  \item \textbf{工具链}：OpenClaw（PM AI）+ Cursor（开发 AI），Claude/Codex 暂不参与
  \item \textbf{链路}：魏子政（甲方）$\to$ OpenClaw（PM）$\to$ Cursor（开发），单线沟通
  \item \textbf{拍板机制}：有歧义的环节，OpenClaw 必须先让魏子政拍板决定方向
\end{itemize}

\section{文档一致性规则}

core/ 文档与 MASTER 不一致时：局部细节以 MASTER 为主；重大流程需魏子政拍板。

\section*{更新记录}

\begin{tabular}{lll}
\toprule
版本 & 日期 & 变更内容 \\
\midrule
v1.0 & 2026-05-06 & 全面重写：对齐 MASTER v1.3（D-001~D-016），更新竞品分析（新增 skills.sh），补充邮箱认证/上下架/安装等已完成功能 \\
v0.2 & 2026-05-06 & 对齐 MASTER v1.0 \\
v0.1 & 2026-05-06 & 初始版本 \\
\bottomrule
\end{tabular}

\end{document}
